% ======================================================================
% UVOD (nenumerisano, rimska numeracija stranica kao i predtekst)
% ======================================================================
\phantomsection
\addcontentsline{toc}{chapter}{Uvod}
\chapter*{Uvod}
\label{ch:uvod}

Savremeni informacioni sistemi suočavaju se sa zahtjevom da istovremeno opsluže dvije fundamentalno suprotstavljene klase opterećenja. Transakcijska obrada podataka (eng.\ \textit{Online Transaction Processing}, OLTP) podrazumijeva kratke, atomarne operacije nad pojedinačnim zapisima, uz stroge zahtjeve konzistentnosti i visoku frekvenciju istovremenih korisnika. Analitička obrada (eng.\ \textit{Online Analytical Processing}, OLAP) počiva, nasuprot tome, na dugotrajnim upitima koji agregiraju milione zapisa i traže visok propusni opseg čitanja, prihvatajući pritom nešto veću latenciju pojedinačnog upita u zamjenu za dubinu analitičkog uvida.

Ova napetost nije isključivo tehničke prirode; ona odražava dublji konflikt između potrebe za ažurnošću transakcijskog stanja i potrebe za analitičkim uvidom u to stanje u realnom vremenu. Konflikt se očituje u nizu domena: platforme s korisnički generiranim sadržajem obrađuju hiljade zapisa u sekundi, a istovremeno moraju napajati analitičke nadzorne ploče; finansijski sistemi izvršavaju transakcije dok modeli rizika konsultuju trenutno stanje portfelja; industrijski nadzorni sistemi prikupljaju senzorske podatke visoke frekvencije dok algoritmi detekcije anomalija zahtijevaju pristup cjelokupnoj historiji mjerenja. U svim tim slučajevima kašnjenje između nastanka transakcijskog zapisa i njegove analitičke dostupnosti neposredno utiče na kvalitetu poslovnih odluka.

Microsoft SQL Server 2022 spada među najraširenije sisteme za upravljanje relacijskim bazama podataka u komercijalnom okruženju, s posebno visokom zastupljenošću u preduzećima srednje i jugoistočne Evrope. Za adresiranje opisanog problema sistem nudi dva ugrađena mehanizma: In-Memory OLTP engine (Hekaton) za ubrzanje transakcija \cite{diaconu2013hekaton} te columnstore indekse za ubrzanje analitike \cite{larson2013columnstore}. Njihova optimalna kombinacija pod mješovitim OLTP/OLAP opterećenjima, međutim, do sada nije sistematski istražena na način koji bi praktičarima ponudio konkretne smjernice.

Aktuelne HTAP arhitekture moguće je razumjeti tek u svjetlu evolucije sistema za upravljanje bazama podataka. Coddov relacijski model iz 1970.\ utemeljio je paradigmu normalizovanih shema optimiziranih za integritet i efikasnost točkačnih operacija \cite{ozcan2017htap}. Krajem 1980-ih i tokom 1990-ih, Inmon (1992.) i Kimball (1996.) formalizirali su koncept skladišta podataka (eng.\ \textit{data warehouse}) kao zasebnog sistema optimiziranog za analitiku. ETL procesi koji su premostili dva sloja postali su industrijski standard, ali i izvor kašnjenja koje se mjeri satima, pa i čitavim danima, između nastanka transakcijskog zapisa i njegove analitičke dostupnosti.

Pad cijene RAM memorije krajem 2000-ih otvorio je prostor za in-memory baze podataka (VoltDB, SAP HANA) koje su eliminirale I/O kao ograničavajući faktor. Microsoft je 2012.\ godine u SQL Server uveo columnstore indekse, a 2014.\ Hekaton in-memory OLTP engine, postavljajući temelje hibridnoj HTAP arhitekturi unutar etabliranog relacijskog sistema. Iste, 2014.\ godine, Gartner je skovao termin HTAP za klasu sistema koji teže eliminaciji ETL-a \cite{ozcan2017htap}. Distribuirani dual-engine sistemi poput TiDB-a \cite{huang2020tidb} i Google F1 Lightning \cite{yang2020f1} potom su demonstrirali izvedivost HTAP-a na petabajt-skali, dok SQL Server 2022 predstavlja aktuelnu razinu zrelosti te evolucije unutar Microsoft ekosistema.

Uprkos teorijskoj obrađenosti HTAP paradigme i dostupnosti praktičnih implementacija, literatura ostaje oskudna studijama koje bi sistematski uporedile različite konfiguracijske strategije pohrane unutar jedne te iste platforme, kakva je SQL Server 2022, i to nad stvarnim skupovima podataka s prirodnom distribucijom aktivnosti. Postojeće benchmarking studije najčešće se okreću distribuiranim NewSQL sistemima (TiDB, SingleStore) ili se oslanjaju isključivo na sintetičke skupove podataka (TPC-C, TPC-H) koji ne odražavaju Pareto distribuciju korisničke aktivnosti svojstvenu web aplikacijama \cite{zhang2024hybench}.

U radu se mogu uočiti tri istraživačka pravca. Prvo, rad donosi empirijsku komparaciju četiri konfiguracijske strategije SQL Servera 2022 nad stvarnim skupom podataka s $\sim$19,2 miliona zapisa. Drugo, za svaku konfiguraciju kvantificira degradaciju OLTP propusnosti pri mješovitom opterećenju i identificira arhitekturalni mehanizam koji stoji iza uočenih razlika. Treće, iz eksperimentalnih nalaza deriviraju se praktične smjernice za odabir konfiguracijske strategije u ovisnosti o omjeru opterećenja i razini konkurentnosti -- smjernice neposredno primjenjive za inženjere i arhitekte koji planiraju HTAP implementaciju u SQL Server okruženju.

Rad je podijeljen na pet poglavlja. Poglavlje~\ref{ch:metodologija} uspostavlja cjeloviti istraživački okvir: iznosi predmet i ciljeve istraživanja, sistematizira teorijsku osnovu HTAP paradigme s pregledom komercijalnih implementacija -- SAP HANA, TiDB, Google F1 Lightning i SQL Server 2022 -- te formalizira tri istraživačke hipoteze uz opis primijenjenog kvantitativnog eksperimentalnog pristupa.

Poglavlje~\ref{ch:sistem} opisuje arhitekturu posmatranog sistema kroz četiri konfiguracijske strategije pohrane -- Baseline, CS, IM i HTAP -- implementirane kao zasebne sheme unutar jedne SQL Server 2022 instance s identičnom logičkom strukturom tabela. Prikazane su ključne arhitekturalne razlike između konfiguracija te karakteristike skupa podataka StackOverflow2010 s $\sim$19,2 miliona zapisa i Pareto distribucijom korisničke aktivnosti.

Poglavlje~\ref{ch:eksperiment} detaljno izlaže eksperimentalni dizajn ($4\times3\times3$ puni faktorijalni, 30 ponavljanja po kombinaciji), hardversko-softversko okruženje i implementaciju .NET 8 benchmark aplikacije, a radna opterećenja definira kroz šest OLTP i šest OLAP pohranjenih procedura. Statističku metodologiju čine Kruskal-Wallis H-test, Mann-Whitney U test s Bonferroni korekcijom, Cohenov $d$ i OLS regresijska analiza s logaritamskom transformacijom zavisnih varijabli.

Poglavlje~\ref{ch:rezultati} prezentira mjerene rezultate za sve četiri konfiguracije, provjerava svaku od tri hipoteze uz statističke dokaze i kvantificirane veličine efekta te interpretira arhitekturalne mehanizme koji stoje iza uočenih razlika. Poglavlje uključuje i usporedbu s dosadašnjom literaturom, konkretne praktične preporuke za odabir konfiguracijske strategije te diskusiju ograničenja validnosti eksperimentalnih zaključaka.

Zaključak sumira ključne nalaze eksperimenta, iznosi originalne istraživačke zaključke te predlaže četiri smjernice za buduća istraživanja, uključujući evaluaciju distribuiranih HTAP arhitektura i uticaja SSD/NVMe pohrane na izmjerene trendove.

Rad uključuje dva dodatka namijenjena reproducibilnosti i transparentnosti istraživanja. Prilog~\ref{app:oltp} sadrži implementaciju reprezentativnih OLTP pohranjenih procedura korištenih u benchmarku, uključujući disk-zasnovane i memorijsko-optimizirane (nativno kompajlirane) varijante, čime se čitatelju pruža uvid u konkretne razlike između ispitivanih konfiguracija na razini izvršnog koda. Prilog~\ref{app:olap} donosi implementaciju reprezentativnih OLAP analitičkih upita primijenjenih u eksperimentu, ilustrirajući složenost i obim analitičkog radnog opterećenja koje je podvrgnuto mjerenju.
